
\documentclass{article}
\usepackage{colortbl}
\usepackage{makecell}
\usepackage{multirow}
\usepackage{supertabular}

\begin{document}

\newcounter{utterance}

\centering \large Interaction Transcript for game `cladder', experiment `full\_v1.5\_default', episode 6328 with qwen3.5{-}9b.
\vspace{24pt}

{ \footnotesize  \setcounter{utterance}{1}
\setlength{\tabcolsep}{0pt}
\begin{supertabular}{c@{$\;$}|p{.15\linewidth}@{}p{.15\linewidth}p{.15\linewidth}p{.15\linewidth}p{.15\linewidth}p{.15\linewidth}}
    \# & \multicolumn{2}{c}{Player} && \multicolumn{2}{c}{Game Master} \\
    \hline

    \theutterance \stepcounter{utterance}  
    & & & \multicolumn{4}{p{0.6\linewidth}}{
        \cellcolor[rgb]{0.9,0.9,0.9}{
            \makecell[{{p{\linewidth}}}]{
                \texttt{\tiny{[P1$\langle$GM]}}
                \texttt{You are an expert in causal inference. The following question is not a typical commonsense query, but rather a meticulously designed question created by a professor specializing in causal inference, intended to assess the students' mastery of the course content.} \\
\texttt{Imagine a self{-}contained, hypothetical world with only the following conditions, and without any unmentioned factors or causal relationships: Swoy has a direct effect on rixq and xevu. Rixq has a direct effect on zuph. Zuph has a direct effect on xevu. Swoy is unobserved. For those who are not rixq and are not zuph, the probability of xevu is 46\%. For those who are not rixq and are zuph, the probability of xevu is 41\%. For those who are rixq and are not zuph, the probability of xevu is 66\%. For those who are rixq and are zuph, the probability of xevu is 60\%. For those who are not rixq, the probability of zuph is 61\%. For those who are rixq, the probability of zuph is 34\%. For those who are rixq, would it be more likely to see xevu if the individual was not rixq?} \\
\texttt{Start your answer with 'Yes' or 'No', followed by additional reasoning or evidence to support your explanation.} \\
            }
        }
    }
     \\ \\

    \theutterance \stepcounter{utterance}  
    & \multicolumn{4}{p{0.6\linewidth}}{
        \cellcolor[rgb]{0.561,0.737,0.561}{
            \makecell[{{p{\linewidth}}}]{
                \texttt{\tiny{[P1$\rangle$GM]}}
                \texttt{No} \\
\\ 
\texttt{**Reasoning:**} \\
\\ 
\texttt{To answer this question, we must compare the probability of observing \$xevu\$ under two specific counterfactual conditions for an individual who is actually \$rixq\$:} \\
\texttt{1.  **Actual State:** The individual is \$rixq\$.} \\
\texttt{2.  **Counterfactual State:** The individual is *not* \$rixq\$ (while keeping other relevant factors constant, specifically the value of \$zuph\$).} \\
\\ 
\texttt{The question asks: "For those who are rixq, would it be more likely to see xevu if the individual was not rixq?"} \\
\\ 
\texttt{Let's break down the probabilities provided in the prompt:} \\
\\ 
\texttt{*   **Probability of \$xevu\$ given \$rixq\$ and \$zuph\$:**} \\
\texttt{    The prompt states: "For those who are rixq and are zuph, the probability of xevu is 60\%."} \\
\texttt{    \$\$P(xevu \textbackslash{}mid rixq, zuph) = 0.60\$\$} \\
\\ 
\texttt{*   **Probability of \$xevu\$ given *not* \$rixq\$ and \$zuph\$:**} \\
\texttt{    The prompt states: "For those who are not rixq and are zuph, the probability of xevu is 41\%."} \\
\texttt{    \$\$P(xevu \textbackslash{}mid \textbackslash{}neg rixq, zuph) = 0.41\$\$} \\
\\ 
\texttt{The causal structure implies that \$rixq\$ is a direct cause of \$zuph\$ (since \$rixq\$ affects \$zuph\$, and \$zuph\$ affects \$xevu\$). However, the question specifically asks us to evaluate the likelihood of \$xevu\$ if the individual were *not* \$rixq\$. In causal inference, when comparing the effect of a treatment (\$rixq\$) to a control (\$\textbackslash{}neg rixq\$), we typically hold the mediator (\$zuph\$) constant to isolate the direct effect, or we look at the total effect by marginalizing over \$zuph\$.} \\
\\ 
\texttt{However, the phrasing "if the individual was not rixq" usually implies a counterfactual comparison where we switch the value of the variable \$rixq\$ from true to false.} \\
\texttt{*   In the actual world (for a \$rixq\$ individual), if they also happen to be \$zuph\$, the probability of \$xevu\$ is **60\%**.} \\
\texttt{*   In the counterfactual world (if that same individual were *not* \$rixq\$), we must determine their status regarding \$zuph\$.} \\
\texttt{    *   If we hold \$zuph\$ constant (comparing \$P(xevu \textbackslash{}mid rixq, zuph)\$ vs \$P(xevu \textbackslash{}mid \textbackslash{}neg rixq, zuph)\$): The probability drops from 60\% to 41\%.} \\
\texttt{    *   If we consider the marginal probability for a non{-}\$rixq\$ individual (averaging over \$zuph\$):} \\
\texttt{        *   \$P(zuph \textbackslash{}mid \textbackslash{}neg rixq) = 0.61\$} \\
\texttt{        *   \$P(\textbackslash{}neg zuph \textbackslash{}mid \textbackslash{}neg rixq) = 1 {-} 0.61 = 0.39\$} \\
\texttt{        *   \$P(xevu \textbackslash{}mid \textbackslash{}neg rixq) = (0.61 \textbackslash{}times 0.41) + (0.39 \textbackslash{}times 0.46) = 0.2501 + 0.1794 = 0.4295\$ (approx 43\%).} \\
\\ 
\texttt{In both interpretations (holding the mediator constant or looking at the marginal probability), the probability of seeing \$xevu\$ is **lower** if the individual is not \$rixq\$ (41\% or \textasciitilde{}43\%) compared to if they are \$rixq\$ (60\% or 66\%).} \\
\\ 
\texttt{Therefore, it would **not** be more likely to see \$xevu\$ if the individual was not \$rixq\$; it would be less likely. The presence of \$rixq\$ increases the probability of \$xevu\$.} \\
            }
        }
    }
    & & \\ \\

    \theutterance \stepcounter{utterance}  
    & & & \multicolumn{2}{p{0.3\linewidth}}{
        \cellcolor[rgb]{0.9,0.9,0.9}{
            \makecell[{{p{\linewidth}}}]{
                \texttt{\tiny{[GM$|$GM]}}
                \texttt{No} \\
            }
        }
    }
    & & \\ \\

    \theutterance \stepcounter{utterance}  
    & & & \multicolumn{2}{p{0.3\linewidth}}{
        \cellcolor[rgb]{0.9,0.9,0.9}{
            \makecell[{{p{\linewidth}}}]{
                \texttt{\tiny{[GM$|$GM]}}
                \texttt{no} \\
            }
        }
    }
    & & \\ \\

    \theutterance \stepcounter{utterance}  
    & & & \multicolumn{2}{p{0.3\linewidth}}{
        \cellcolor[rgb]{0.9,0.9,0.9}{
            \makecell[{{p{\linewidth}}}]{
                \texttt{\tiny{[GM$|$GM]}}
                \texttt{game\_result = WIN} \\
            }
        }
    }
    & & \\ \\

\end{supertabular}
}

\end{document}
